% Generated draft: requires mathematical and source review.
\documentclass[11pt]{article}
\usepackage[margin=25mm]{geometry}
\usepackage{fontspec}
\setmainfont{Latin Modern Roman}
\usepackage{amsmath,amssymb,mathtools,unicode-math}
\setmathfont{Latin Modern Math}
\usepackage{enumitem,needspace,xurl,hyperref}
\setlength{\parindent}{0pt}
\setlength{\parskip}{5pt}
\setlength{\emergencystretch}{3em}
\newcommand{\N}{\mathbb{N}}
\newcommand{\Z}{\mathbb{Z}}
\newcommand{\Q}{\mathbb{Q}}
\newcommand{\R}{\mathbb{R}}
\newcommand{\C}{\mathbb{C}}
\newcommand{\F}{\mathbb{F}}
\newcommand{\A}{\mathbb{A}}
\newcommand{\PP}{\mathbb{P}}
\newcommand{\E}{\mathbb{E}}
\newcommand{\eps}{\varepsilon}
\newcommand{\dd}{\,\mathrm{d}}
\begin{document}
\section*{UnsolvedMath ID: 30000205}
\textbf{Source problem code:} catalogue entry
\subsection{Definitions and mathematical statement}
\textbf{Tits' Center Conjecture, as supplied.} Let \(X\) be a spherical building and \(Y\) a nonempty convex contractible subcomplex. Must there exist a point \(y\in|Y|\) such that\[g(y)=y\quad\text{for every }g\in\operatorname{Aut}(X)\text{ with }g(Y)=Y?\] Here \(|Y|\) is the geometric realization, \(g(Y)=Y\) means setwise stabilization, and the quantifier asks for one common fixed point of the entire stabilizer. Separate fixed points chosen for individual automorphisms would be a weaker conclusion. The fixed point is required to lie in \(Y\), not merely somewhere in \(X\).\par In the usual simplicial convention, a spherical building is a building whose apartments are spherical Coxeter complexes, with its standard piecewise spherical metric. Convexity of \(Y\) means that each geodesic segment of length less than \(\pi\) between points of \(|Y|\) lies in \(|Y|\). Contractibility refers to the topology of that realization. Automorphisms are the intended simplicial building automorphisms; the source does not state a restriction to type-preserving ones, so that restriction is not inserted silently. It asks for a point, not necessarily a vertex. The excerpt supplies no thickness, local-finiteness, rank, or building-type restrictions, and any such hypotheses in the intended source version must be checked.\par The supplied triage checked 2026-08-21 labels the record open solely while describing its extraction as a theorem-followup fragment under a triage rule. The original and curated fields now state the center conjecture explicitly and the verification field reports checking the official source. The fragment-based status rationale does not assess that complete assertion or later results. Accordingly, the historical statement and Tits attribution are retained without asserting that the conjecture remains unresolved today.
\subsection{Short English statement}
A convex contractible subcomplex should contain one point fixed by every building automorphism preserving it. The stabilizer acts setwise, and a common fixed point is required. The supplied open-status rationale describes an incomplete fragment rather than assessing the full conjecture now displayed.
\subsection{Sources}
\begin{itemize}\item Local UnsolvedMath registry, permanent ID 30000205, code catalogue entry; \emph{Tits' Center Conjecture for Spherical Buildings}, attributed to Tits in the original statement; Oberwolfach Reports Open Problems.\item \emph{Groups and Geometries (2005)}, Oberwolfach Reports, DOI 10.4171/owr/2005/12, pp. 647--721. \url{https://doi.org/10.4171/owr/2005/12}. Proposed year 2005; no separate proposer field supplied.\item Supplied literature label \emph{Original OWR report}: \url{https://doi.org/10.4171/owr/2005/12}. No separate authoritative original citation or later theorem reference is supplied.\item Literature triage checked 2026-08-21, registry status \texttt{open}, with a fragment-based triage qualification; record updated 2026-09-16. The \texttt{corrected\_verified} field reports official-report inspection, not independently checked here.\end{itemize}
\end{document}
